\chapter{\label{chap:chap5}Recursos Utilizados}

O desenvolvimento do \textbf{Dashboard Financeiro} utilizou ferramentas modernas e acessíveis, adequadas ao desenvolvimento web full-stack, além de recursos computacionais que garantiram eficiência e organização ao longo do projeto.

\section{Recursos Computacionais}

O projeto foi desenvolvido em um \textbf{computador pessoal} capaz de executar simultaneamente servidores locais, contêineres Docker e editores de código sem degradação perceptível.

Houve acesso contínuo à \textbf{internet}, usado para consultar documentações técnicas, integrar bibliotecas externas, acessar plataformas de versionamento e manter os repositórios sincronizados.


\section{Modelo de Domínio}

Antes de derivar o banco de dados, o domínio foi modelado considerando entidades, agregados e relacionamentos principais (Figura~\ref{fig:modelo-dominio}). O diagrama foi gerado com PlantUML a partir das entidades reais do backend e segue a notação de classes UML para evidenciar multiplicidades e atributos relevantes às regras de negócio.

\begin{figure}[H]
\centering
\safeincludegraphics[width=0.95\textwidth]{figs/modelo-dominio.png}
\caption{Modelo de domínio do \textit{Dashboard Financeiro}.}
\label{fig:modelo-dominio}
\end{figure}

Os detalhes do modelo relacional e dos esquemas versionados por Liquibase foram consolidados na Seção~\ref{sec:banco-dados}, garantindo que a leitura sobre arquitetura (Capítulo~\ref{chap:chap3}) siga exatamente a ordem personas $\rightarrow$ histórias de usuário $\rightarrow$ arquitetura $\rightarrow$ banco de dados.

\section{Implantação em Nuvem, Domínio e E-mails}

O Volume I previa o deploy no Google Cloud, porém durante o TC~II optou-se pelo \textbf{Railway}, que oferece \textit{build} automatizado a partir dos Dockerfiles já presentes no repositório e gerenciamento simplificado de múltiplos serviços com variáveis de ambiente compartilhadas. Cada microserviço (cadastro/autenticação, financeiro, notificações, gateway e frontend) possui imagem própria publicada no Railway e conectada a instâncias gerenciadas de MySQL e Kafka.

Para garantir identidade própria, foi adquirido junto à \textbf{KingHost} o domínio \texttt{tcc\-leonardopreczevski.com.br}. O subdomínio \texttt{app.tcc\-leonardopreczevski.com.br} aponta via CNAME para o frontend publicado no Railway, tornando-se o endereço público adotado em toda a comunicação (frontend, gateway e notificações).

O mesmo domínio foi validado na plataforma \textbf{Resend} para envio de e-mails transacionais. Os registros DNS (SPF, DKIM e \textit{tracking}) foram configurados no painel da KingHost e replicados dentro da Resend, permitindo que o serviço de notificações envie mensagens autenticadas com remetente \texttt{no-reply@app.tcc\-leonardopreczevski.com.br}. Além disso, a Resend utiliza o domínio como assinatura criptográfica, aumentando a entregabilidade dos alertas de metas e das confirmações de cadastro.

Por fim, o gateway e os demais serviços hospedados no Railway reutilizam esse mesmo subdomínio autorizado, garantindo que todas as camadas reconheçam \texttt{app.tcc\-leonardopreczevski.com.br} como origem oficial, em conformidade com o requisito institucional de implantação em nuvem.

\section{Ambiente de Desenvolvimento}

\paragraph{}As principais ferramentas efetivamente adotadas no projeto foram:

\begin{itemize}
    \item \textbf{Visual Studio Code} – IDE utilizada para o desenvolvimento do frontend em React.js \cite{REACT2025};
    \item \textbf{IntelliJ IDEA} – IDE para o desenvolvimento do backend utilizando Spring Boot;
    \item \textbf{Docker} – para criação de contêineres com o banco de dados MySQL, Kafka e execução local dos microserviços;
    \item \textbf{Git e GitHub} – para versionamento e hospedagem do código-fonte, promovendo rastreabilidade, controle de alterações e colaboração;
    \item \textbf{Navegadores modernos} – como Google Chrome e Firefox, para testes de interface e comportamento da aplicação;
    \item \textbf{KingHost} – provedor de domínio e DNS utilizado para validar o envio de e-mails transacionais e preparar a camada pública do sistema;
    \item \textbf{MySQL} – banco de dados relacional executado localmente via contêiner Docker \cite{MYSQL2025};
    \item \textbf{Spring Boot, Spring Security e Spring Cloud Gateway} – frameworks utilizados para estruturação da API RESTful, autenticação, roteamento e segurança \cite{SPRINGSECURITY2025};
    \item \textbf{Apache Kafka} – barramento de eventos para desacoplar microserviços e acionar o serviço de notificações;
    \item \textbf{Resend} – plataforma de e-mails transacionais utilizada em conjunto com domínio próprio.
\end{itemize}

\chapter{\label{chap:chap6}Avaliação da Solução}

Esta seção ficará reservada para consolidar os resultados da pesquisa de satisfação que será aplicada aos usuários do \textbf{Dashboard Financeiro}. O espaço a seguir será preenchido com:
\begin{itemize}
    \item o desenho metodológico da pesquisa (perfil dos respondentes, perguntas e instrumento aplicado);
    \item os indicadores coletados (média de satisfação, NPS, sugestões textuais);
    \item a análise crítica dos dados e potenciais ajustes na solução.
\end{itemize}

Assim que a pesquisa for realizada, os dados serão inseridos aqui antes das Considerações Finais, garantindo que o volume traga evidências empíricas do impacto da aplicação.
